\section{Neural Compute、Memory 与 Personalization}

讲者随后用计算机系统类比 agent：模型像可调用的 neural compute unit，重复调用形成循环，scratchpad 保存工作状态，长期记忆承担跨会话持久化。这个类比很有教学价值，但必须明确：Transformer 不是 CPU，embedding store 也不是磁盘；类比的用途是帮助拆分计算、工作记忆与持久状态。

\subsection{模型作为 Neural Compute Unit}

第一张图把最大输入 token 序列映射到最大输出 token 序列，强调一次模型调用是有限窗口内的计算。第二张 MIPS 指令类比提示：传统处理器有固定指令语义，而语言模型的“指令”由自然语言和上下文共同解释，因此灵活但不稳定。

\lecturefigure{slide-46-neural-compute-unit.jpg}{一次模型调用被抽象成 neural compute unit}{00:28:00}

\lecturefigure{slide-47-mips-instruction-analogy.jpg}{MIPS 固定指令与自然语言模型调用的类比}{00:28:34}

\begin{knowledgebox}{类比中的对应关系}
输入 tokens 类似程序与数据，模型参数类似固定计算结构，输出 tokens 类似计算结果，scratchpad 类似工作区。但语言模型没有硬件 ISA 的严格语义、事务和确定性；因此每次调用都需要 schema、验证和错误处理来补足。
\end{knowledgebox}

\subsection{Looped Transformer：把一次生成变成迭代计算}

Looped Transformer 图把输出重新送入输入，并在 scratchpad、memory、instructions 之间循环。Agent 系统正是用 orchestration 把一次前向生成扩展成多步算法。若第 $k$ 轮内部状态为 $z_k$，可抽象为
\[
z_{k+1}=F_\theta(z_k,o_k,m_k),\qquad
a_k=G(z_{k+1}),
\]
$F_\theta$ 是模型计算，$o_k$ 是新观察，$m_k$ 是检索记忆，$G$ 把内部状态映射到工具动作。系统必须定义何时停止，否则循环本身会成为故障模式。

\lecturefigure{slide-48-looped-transformer.jpg}{Looped Transformer：scratchpad、memory 与 instructions 进入迭代回路}{00:28:51}

\begin{importantbox}{停止条件比“再思考一下”重要}
终止条件至少包括任务完成、预算耗尽、状态重复、风险阈值、用户取消和人工接管。没有停止条件的 reasoning loop 会把不确定性转化为 token、延迟和工具费用。
\end{importantbox}

\subsection{Long-term Memory：持久、检索、更新与遗忘}

长期记忆 slide 把 memory 描述为 long-lived、persistent，并列出 embedding、retrieval models、层级、时间一致性、结构与在线适应。首次出现时必须区分：\term{working memory} 是当前任务状态；\term{long-term memory} 是跨会话持久信息；\term{model memory} 则可能指参数中学到的知识，三者的更新机制不同。

\lecturefigure{slide-49-long-term-memory.jpg}{长期记忆的机制与开放问题}{00:35:17}

一个基础检索式为
\[
m_t=\operatorname{TopK}_{i}\;\operatorname{sim}(q_t,e_i),
\]
$q_t$ 是当前查询 embedding，$e_i$ 是记忆条目表示，$m_t$ 是返回的 top-$k$ 条目。仅按相似度会召回过期或敏感内容，因此实际分数应加入时间、来源、权限与置信度：
\[
\text{score}_i=\alpha\,\text{sim}(q_t,e_i)+\beta\,\text{recency}_i+
\eta\,\text{trust}_i-\lambda\,\text{privacyRisk}_i.
\]

\begin{warningbox}{记忆污染与错误持久化}
一次聊天误解若被写入长期记忆，会在未来被反复强化。写入前应区分用户明确陈述、系统推断和临时任务状态；高影响偏好应允许用户查看、更正和删除。
\end{warningbox}

\begin{lstlisting}[caption={带来源和过期策略的记忆记录}]
memory_item = {
  "fact": normalized_content,
  "source": source_event,
  "confidence": confidence,
  "scope": ["travel", "restaurants"],
  "created_at": timestamp,
  "expires_at": expiry,
  "sensitivity": sensitivity,
  "user_editable": True
}
\end{lstlisting}

\subsection{Personalization 是用户--Agent 对齐问题}

个性化 slide 同时列出显式偏好与隐式偏好。显式偏好如过敏、座位和收藏菜品，通常可以直接确认；隐式偏好来自行为统计，更容易把偶然选择、价格约束或历史环境误当成稳定喜好。系统目标可写为约束优化：
\[
\max_\pi\;\mathbb{E}[U_{\text{user}}(\tau)]
\quad\text{s.t.}\quad
C_{\text{privacy}}(\tau)\le \epsilon_p,
\;C_{\text{safety}}(\tau)\le \epsilon_s.
\]
$U_{\text{user}}$ 是用户效用，$C_{\text{privacy}}$ 与 $C_{\text{safety}}$ 是隐私和安全成本，$\epsilon_p,\epsilon_s$ 是允许阈值。个性化不能以越界收集数据为代价。

\lecturefigure{slide-50-personalization.jpg}{个性化：显式与隐式偏好共同构成 user-agent alignment}{00:36:42}

\lecturefigure{slide-51-personalization-challenges.jpg}{个性化挑战：数据收集、学习、在线适应与隐私}{00:37:56}

两张个性化图应和前面的长期记忆一起阅读：memory 决定“保存什么”，personalization 决定“如何使用”，privacy 决定“什么时候不能用”。例如用户曾为一次商务旅行选择靠走道座位，不应被永久归纳成全局偏好；系统需要保存来源、场景和置信度，并在高影响决策前重新确认。真正的个性化不是让 agent 猜得更多，而是让它知道哪些偏好可信、哪些过期、哪些必须再次询问。

\teachervoice{00:35:17--00:38:00，讲者把长期记忆类比为磁盘，把个性化定义为 user-agent alignment，并强调主动询问、被动学习、SFT、human feedback、on-the-fly adaptation 与 privacy。课堂提示：学习偏好必须带来源和可撤销机制。}

\begin{knowledgebox}{术语消化：Memory 与 Personalization}
\begin{tabularx}{\textwidth}{>{\bfseries}p{0.2\textwidth} X X}
术语 & 核心机制 & 主要失败 \\
\hline
Embedding memory & 相似度检索文本或事件 & 语义近但事实错、时间过期 \\
Episodic memory & 保存具体交互与结果 & 噪声多、隐私敏感 \\
Semantic memory & 把多次事件压缩为稳定事实 & 归纳错误被长期固化 \\
Online adaptation & 根据新反馈即时改变行为 & 灾难性漂移、被恶意反馈操纵 \\
Preference model & 预测用户选择或满意度 & 把历史选择当成真实价值 \\
\end{tabularx}
\end{knowledgebox}

\subsection{本章小结}

Neural-compute 类比帮助把模型调用、循环、工作状态和长期记忆分层，但不能替代真实接口设计。长期记忆需要检索、来源、过期、权限和用户编辑；个性化需要在效用、隐私与安全约束下学习，而不是无限收集行为数据。

